\documentclass[11pt]{article}
\usepackage[T1]{fontenc}
\usepackage[utf8]{inputenc}
\usepackage{lmodern}
\usepackage{amsmath,amssymb,mathrsfs,booktabs,hyperref}
\usepackage[margin=23mm]{geometry}
\hypersetup{hidelinks,pdftitle={Yang--Mills: complete research proofs}}
\setlength{\parindent}{0pt}
\setlength{\parskip}{6pt}
\begin{document}
\section*{The higher carrier and the source-correction evolution}

30 September 2026. This note constructs the missing base-space morphism,
extends the profile map to that carrier, and advances the source equation.
The globally smooth interacting gapless quantum state remains the programme
goal. The results here are classical bundle and evolution results.

\subsection*{1. Exact inputs and conventions}

Retain \(X\cong S^6\), the right principal bundle \(P_H\to X\),

\[
 K=\rho(\operatorname{Sp}(1))\subset G=\operatorname{Spin}(8)
 \subset L=F_4,
 \qquad Q_H=P_H\times_\rho G,
\]

and the source's complete ordered representation

\[
 W=W_r\oplus W_q\oplus W_p,
 \qquad W_r\cong\mathbb R\oplus\operatorname{Im}\mathbb H\oplus\mathbb H,
 \quad W_q\cong\mathbb H\oplus\mathbb H,
 \quad W_p\cong\mathbb H\oplus\mathbb H.
\]

The first action is \(B_u(a,b)=(ua\bar u,b)\). The other two original
upper-entry actions are \(D_u(a,b)=(a\bar u,b\bar u)\); the ordered
conjugations \((a,b)\mapsto(\bar a,\bar b)\) identify each with two
left quaternionic lines. Associated bundles use \([pu,v]=[p,\rho(u)v]\).
The map in \href{https://github.com/KokunoYumeto/yang-mills-interacting-workbench/blob/main/yang-mills/continuations/20260930-s6-ns-moment-map-bridge/PROOF.md}{PROOF.md, Theorem 7.1} is retained literally:

\[
 \mathcal M_H(\zeta;\alpha,\beta;\gamma,\delta)
 =\mu_{\mathbf i}(\alpha)w^{(1)}
  +\mu_{\mathbf j}(\beta)w^{(2)}
  +\mu_{\mathbf k}(\gamma)w^{(3)},\qquad
 \mu_e(a)=ae\bar a.
 \tag{HC1}
\]

Here \(w^{(1)}=\nabla\times(0,0,\chi x_2)\),
\(w^{(2)}=\nabla\times(0,0,-\chi x_1)\), and
\(w^{(3)}=\nabla\times(0,\chi x_1,0)\), with the same
\(0<r<R\), \(\chi=1\) on \(B_r(0)\), and
\(\operatorname{supp}\chi\subset B_R(0)\). No spatial cutoff is removed
from this original map.

The quaternionic bracket is \([v,w]=2v\times w\). In the fundamental
\(SU(2)\) representation,

\[
 T_a=-i\sigma_a/2,\quad [T_a,T_b]=\epsilon_{abc}T_c,\quad
 -2\operatorname{tr}(T_aT_b)=\delta_{ab},\quad
 \varphi(\mathbf i,\mathbf j,\mathbf k)=2(T_1,T_2,T_3).
 \tag{HC2}
\]

The physical metric is \(\eta=\operatorname{diag}(-1,1,1,1)\),
\(s=x^0=ct\), \(c>0\), and \(D_\mu=\partial_\mu+[A_\mu,\,\cdot\,]\).
In particular \(\partial_s=c^{-1}\partial_t\). The source convention is
\(D^\mu F_{\mu\nu}=g^2J_\nu\), \(g>0\).

The exact retained source is
\href{https://github.com/KokunoYumeto/yang-mills-interacting-workbench/blob/main/yang-mills/continuations/20260930-s6-ns-moment-map-bridge/sources/higher_rung/s6_higher_rung_24d_preprint.tex}{the higher-rung TeX},
SHA-256 `8c526746de9b56a4fcd0a274df0c470092cc82a16e49857df539a1c9d033956f`.
Its lines 2593–2804 prove the tangent comparison, and lines 5645–5792
prove that \(Q_H\) is trivial while its specified reduction remains
nontrivial. Those are inputs, with their full source proofs retained.

\subsection*{2. The reduction carrier and its exact pullback}

Define

\[
 M=L/G,\qquad D=L/K,\qquad
 \pi:D\longrightarrow M,\quad \ell K\longmapsto\ell G,
 \qquad j:G/K\longrightarrow D,\quad aK\longmapsto aK.
 \tag{HC3}
\]

The embedding of \(K\) into \(L\) is precisely the source embedding through
the ordered Jordan-frame stabilizer. It is not an independently chosen
copy of \(\operatorname{Sp}(1)\).

Theorem 2.1. The map \(\pi\) is a smooth fibre bundle with fibre \(G/K\).
The space \(D\) has dimension \(52-3=49\), the base \(M\) has dimension
\(52-28=24\), and each fibre has dimension \(28-3=25\). There is a map
\(r_H:X\to D\) such that

\[
 \pi r_H(x)=G,\qquad
 r_H^*(L\to D)\cong P_H,
 \qquad r_H^*(L\times_K W)\cong\mathcal W_H.
 \tag{HC4}
\]

The first pullback uses \(\rho\) to identify the structure groups.
Both \(r_H\) and its fibre map \(f_H:X\to G/K\) are non-null-homotopic.

Proof. A local section \(\sigma:U\to L\) of \(L\to L/G\) gives
the explicit trivialization

\[
 U\times G/K\longrightarrow\pi^{-1}(U),\qquad
 (y,aK)\longmapsto\sigma(y)aK.
\]

Its inverse sends \(\ell K\) to
\((\ell G,\sigma(\ell G)^{-1}\ell K)\). On an overlap
\(\sigma'=\sigma h\), the fibre coordinate changes by \(aK\mapsto h^{-1}aK\).
These formulas prove local triviality and smoothness. The group dimensions
give the stated dimensions. The fibre over \(G\) is exactly \(j(G/K)\).

Choose a principal-bundle trivialization
\(\tau:Q_H\to X\times G\). For \(p\in(P_H)_x\), write
\(\tau([p,1])=(x,a_p)\). The retained associated-bundle convention gives
\(a_{pu}=a_p\rho(u)\). Thus \(f_H(x)=a_pK\) is independent of \(p\),
and \(r_H=jf_H\) is well defined. The explicit pullback isomorphism is

\[
 P_H\longrightarrow r_H^*(L\to D),\qquad
 p\longmapsto(x,a_p).
 \tag{HC5}
\]

Every element \((x,\ell)\) of the target has \(\ell K=a_pK\), hence
\(\ell=a_p k\) for a unique \(k\in K\). There is a unique \(u\in\operatorname{Sp}(1)\)
with \(k=\rho(u)\), and the inverse of (HC5) is \(pu\). This proves
bijectivity, equivariance, and continuity, and proves smoothness whenever
smooth bundle structures and \(\tau\) have been chosen. Applying the
representation \(W\) gives the last isomorphism in (HC4).

If \(r_H\) were null-homotopic, its pullback principal bundle would be
isomorphic to the pullback along a constant map, which is trivial. This
contradicts the retained nonzero order-two clutching class of \(P_H\).
The same proof applies to \(f_H\). Also \(\pi r_H\) is constant by (HC3).
This proves every assertion. \(\square\)

The dependence on \(\tau\) is explicit. If
\(\tau'=b(x)\tau\), with \(b:X\to G\), then
\(f'_H(x)=b(x)f_H(x)\) and \(r'_H(x)=j(b(x)f_H(x))\). No canonical
trivialization, embedding of \(X\), or reduction of \(TM\) over \(M\)
is being inferred. The exact relation is the pullback (HC4).

Smooth-structure detail. The literal source map defining \(P_H\) is
retained as a continuous map; its differentiability at quotient seams is
not needed. A continuous clutching map \(S^5\to\operatorname{Sp}(1)=S^3\)
can be uniformly approximated in \(\mathbb R^4\) by a smooth map within
distance \(1/2\). Radial projection of that approximation onto \(S^3\)
is smooth. Radially projecting the straight homotopy gives a homotopy
of clutching maps, because every intermediate vector has norm at least
\(1/2\). The resulting smooth bundle is isomorphic to the original
topological bundle: a homotopy of clutching maps glues the product disks
over \([0,1]\), and local trivializations along finitely many subintervals
give the endpoint isomorphism. Transport its smooth structure through
that isomorphism. The null-homotopy of the extended clutching map into
\(G\subset O(N)\) can likewise be smoothed relative to its boundary:
smooth its matrix entries with a collar held fixed, then apply the smooth
nearest-point projection of a tubular neighbourhood of the compact
embedded group \(G\). This gives a smooth trivialization \(\tau\).
Accordingly \(r_H\) may be chosen smooth, without replacing the source
map \(H\) or claiming that its original seams were smooth.

\subsection*{3. The rank-24 carrier, with the spectral factors retained}

Put

\[
 E_D=L\times_K W\longrightarrow D,\qquad
 \mathcal A_D=L\times_K\operatorname{Im}\mathbb H.
\]

Choose the original distinct real parameters
\(\lambda_1+\lambda_2+\lambda_3=0\), and retain

\[
 d_\lambda=\sum_i\lambda_i e_i,\quad
 \Delta_1=\lambda_2-\lambda_3,\quad
 \Delta_2=\lambda_3-\lambda_1,\quad
 \Delta_3=\lambda_1-\lambda_2.
 \tag{HC6}
\]

The orbit comparison is \(\Xi_\lambda:M\to Ld_\lambda\),
\(\ell G\mapsto\ell d_\lambda\). For \(w=(r,q,p)\), let
\(J_0(w)=J(0;p,q,r)\) be the original zero-diagonal Albert matrix.

Proposition 3.1. The formula

\[
 \Psi:E_D\longrightarrow\pi^*TM,\qquad
 [\ell,w]\longmapsto
 \left(\ell K,(d\Xi_\lambda)^{-1}_{\ell G}
                         (\ell J_0(w))\right)
 \tag{HC7}
\]

is a bundle isomorphism. All three eight-dimensional labels are retained.

Proof. The off-diagonal Albert matrices are exactly the tangent space
at \(d_\lambda\): in cyclic source coordinates the three tangent maps obey
\(D_i(a)d_\lambda=\Delta_iE_i(a)\), and each \(\Delta_i\ne0\).
The source proves that these 24 directions exhaust the tangent space.
For \(k\in K\), \(J_0(kw)=kJ_0(w)\); hence replacing
\((\ell k,w)\) by \((\ell,kw)\) leaves (HC7) unchanged. The map is linear
and invertible in every fibre. Local trivializations make the inverse
smooth. In a cyclic slot \(E_i(a)\), the inverse tangent coordinate is
\(K_i(a/\Delta_i)\), exactly as in the retained source. Thus the trace
metric gives \(2h(a,b)\) on those original orbit coordinates, while
\(K_i(a)\) gives \(2\Delta_i^2h(a,b)\). Neither the sign nor the scale
of a \(\Delta_i\) has been removed. \(\square\)

The exact sequence of tangent bundles of \(D\to M\) has vertical part
\(L\times_K(\mathfrak g/\mathfrak k)\). At the Lie-algebra level it is

\[
 0\longrightarrow\mathfrak g/\mathfrak k
 \longrightarrow\mathfrak l/\mathfrak k
 \longrightarrow\mathfrak l/\mathfrak g\longrightarrow0.
 \tag{HC8}
\]

Its dimensions are \(25,49,24\). The rank-24 bundle \(E_D\) is the
pullback of \(TM\); it is not all of \(TD\). Formula (HC8), with (HC7),
proves the precise relation among those objects.

The principal extension has an equally explicit comparison:

\[
 L\times_KG\longrightarrow\pi^*(L\to M),\qquad
 [\ell,a]\longmapsto(\ell K,\ell a).
\]

It is well defined because \([\ell k,a]=[\ell,ka]\) gives the same
second coordinate. Given \((d,b)\) with \(bG=\pi(d)\), choose a
representative \(\ell\) of \(d\). Then \(a=\ell^{-1}b\in G\), and the
inverse is \([\ell,\ell^{-1}b]\). Replacing \(\ell\) by \(\ell k\)
gives the same class. The formulas and local sections prove that these
are smooth inverse principal-bundle morphisms. Pullback by \(r_H\)
therefore recovers the source's trivial extension \(Q_H\), together
with its specified nontrivial \(K\)-reduction.

The same polynomial (HC1) now defines

\[
 \mathcal M_D:E_D\longrightarrow
 \mathcal A_D\otimes\underline{\mathcal V_{\mathrm{div},R}}.
 \tag{HC9}
\]

Its equivariance is the identity
\(\mu_e(ua)=u\mu_e(a)\bar u\). Therefore it is well defined on associated
classes. Pullback by \(r_H\) gives exactly \(\mathcal M_H\), since (HC5)
carries each labelled coordinate and each fixed \(w^{(a)}\) unchanged.
The proof of Theorem 7.1 applies in every fibre: the zero subbundle has
rank 12 and the vertical ranks are \(0,3,6,9\). This proves extension to
the higher carrier, including its commuting pullback square.

\subsection*{4. A full connection on the parameter space times spacetime}

A fibre input is not a chosen global section of \(E_D\). To retain every
input without imposing such a section, use the total space \(Y=\operatorname{Tot}(E_D)\)
with projection \(p:Y\to D\). The pullback \(p^*E_D\) has the tautological
section \(v\mapsto(v,v)\). Applying (HC9) to it defines a smooth adjoint-valued
spatial one-form \(C\) on \(Y\times\mathbb R^3\).

Choose the unrescaled trace inner product
\(b(U,V)=-\operatorname{Tr}_{J}(UV)\) on \(\mathfrak l\subset\operatorname{End}(J)\).
Automorphisms act orthogonally for the Albert trace metric, so the
derivations are skew-adjoint. Thus \(b(U,U)>0\) for \(U\ne0\), and cyclicity
of trace proves \(\operatorname{Ad}(L)\)-invariance. Let
\(\mathfrak n=\mathfrak k^{\perp_b}\) and define

\[
 \omega_K=\operatorname{pr}_{\mathfrak k}(\ell^{-1}d\ell),\qquad
 \theta=\operatorname{pr}_{\mathfrak n}(\ell^{-1}d\ell).
 \tag{HC10}
\]

The projection commutes with \(\operatorname{Ad}(K)\). Therefore
\(R_k^*\omega_K=\operatorname{Ad}_{k^{-1}}\omega_K\), and \(\omega_K\)
returns the generator on every fundamental vertical vector. These two
identities prove that it is a principal connection on \(L\to D\).
The Maurer–Cartan identity gives its complete curvature

\[
 \Omega_K=-\tfrac12\operatorname{pr}_{\mathfrak k}[\theta\wedge\theta].
 \tag{HC11}
\]

Indeed \([\mathfrak k,\mathfrak n]\subset\mathfrak n\), so the projected
mixed bracket vanishes, and the \([\omega_K\wedge\omega_K]/2\) terms
cancel exactly. No assertion of flatness follows.

Pull this connection to \(Y\), then use \((d\rho)^{-1}\) and \(\varphi\)
to express it in the same \(SU(2)\) convention as (HC2). Call the resulting
local base form \(b_A\,dy^A\). There is now a full connection

\[
 \mathbb B=b_A\,dy^A+\varphi(C_i)\,dx^i,
 \qquad \mathbb B_0=0,
 \tag{HC12}
\]

on the pulled-back principal bundle over \(Y\times\mathbb R^{1,3}\).
On base overlaps with gauge change \(u(y)\), the spatial field changes by
adjoint conjugation and the base form acquires \(u^{-1}du\). Their sum
therefore obeys the full connection law. Its curvature components are

\[
 \begin{aligned}
 \mathbb F_{AB}&=\partial_A b_B-\partial_B b_A+[b_A,b_B],\\
 \mathbb F_{Ai}&=\partial_A\varphi(C_i)+[b_A,\varphi(C_i)],\\
 \mathbb F_{ij}&=\partial_i\varphi(C_j)-\partial_j\varphi(C_i)
                       +[\varphi(C_i),\varphi(C_j)],\\
 \mathbb F_{A0}&=0,\qquad\mathbb F_{0i}=0.
 \end{aligned}\tag{HC13}
\]

The base curvature persists outside the spatial support of \(C\); it is
explicitly retained. The fibre restriction to a fixed \(y\) is the
four-dimensional connection studied below. No metric or Yang–Mills
equation in the 73 parameter dimensions of \(Y\) is imposed by that
four-dimensional restriction. Imposing such an equation would require
additional metric and dynamical data. Formula (HC13) exhibits the mixed
terms that must be kept when parameter derivatives are used.

\subsection*{5. The exact correction equation and initial constraint}

Work in the quaternionic convention until applying (HC2). For a fixed
fibre input, write \(C_i(x)=(\mathcal M_H(v))_i(x)\), \(C_0=0\), and
define all the following objects with spatial indices \(i,j=1,2,3\):

\[
 F_{ij}=\partial_iC_j-\partial_jC_i+[C_i,C_j],\quad
 D_i=\partial_i+[C_i,\,\cdot\,],\quad
 S_j=\sum_iD_iF_{ij}.
 \tag{HC14}
\]

For a temporal-gauge correction \(a_i(s,x)\), put

\[
 K_{ij}=D_i a_j-D_j a_i,\quad Q_{ij}=[a_i,a_j],\quad
 F_{ij}(C+a)=F_{ij}+K_{ij}+Q_{ij}.
\]

The source-free equations are exactly

\[
 \begin{aligned}
 0={}&\sum_i\left(D_i\dot a_i+[a_i,\dot a_i]\right),\\
 \ddot a_j={}&S_j+\sum_i\bigl(
 D_iK_{ij}+D_iQ_{ij}+[a_i,F_{ij}]
                  +[a_i,K_{ij}]+[a_i,Q_{ij}]\bigr).
 \end{aligned}\tag{HC15}
\]

A dot means \(\partial_s\), so \(\ddot a=c^{-2}\partial_t^2a\).
To prove (HC15), use \(F_{0j}=\dot a_j\), \(F_{i0}=-\dot a_i\),
\(D^0=-\partial_s\), and \(D_i^{C+a}=D_i+[a_i,\cdot]\), and expand
every term in \(D^\mu F_{\mu\nu}(C+a)\). There are no omitted powers of
\(a\). This derivation proves (HC15) for all fibre inputs, including every
rank stratum of (HC1).

In particular \(a(0)=0\), \(\dot a(0)=0\) satisfies the Gauss equation
for every original compactly supported input. The evolution equation
requires \(\ddot a_j(0)=S_j\). On the three-colour core this is
\(-8(\mathbf i,\mathbf j,\mathbf k)\), or
\(-16(T_1,T_2,T_3)\) after applying \(\varphi\). The static source is
therefore a prescribed initial acceleration, not a failure of the
initial Gauss constraint.

\subsection*{6. An exact Gauss-preserving compactly supported correction}

Proposition 6.1. The field

\[
 A_i(s,x)=C_i(x)+\frac{s^2}{2}S_i(x),\qquad A_0=0
 \tag{HC16}
\]

is smooth for all real \(s\), is supported in the same spatial compact
set as \(C\), has finite spatial energy at each \(s\), satisfies Gauss's
law exactly for every \(s\), and has spatial source

\[
 D_A^\mu F_{\mu j}(A)
 =\frac{s^2}{2}L_j+\frac{s^4}{4}Q_j+\frac{s^6}{8}N_j,
 \tag{HC17}
\]

where the complete coefficients are

\[
 \begin{aligned}
 B_{ij}&=D_iS_j-D_jS_i,\qquad C_{ij}=[S_i,S_j],\\
 L_j&=\sum_i\left(D_iB_{ij}+[S_i,F_{ij}]\right),\\
 Q_j&=\sum_i\left(D_iC_{ij}+[S_i,B_{ij}]\right),\\
 N_j&=\sum_i[S_i,C_{ij}].
 \end{aligned}\tag{HC18}
\]

The original three-colour core remains non-Abelian for all sufficiently
small \(|s|\).

Proof. The spatial covariant-divergence identity is

\[
 2\sum_jD_jS_j
 =\sum_{i,j}(D_jD_i-D_iD_j)F_{ij}
 =\sum_{i,j}[F_{ji},F_{ij}]=0.
 \tag{HC19}
\]

The first equality follows by exchanging \(i,j\) and using
\(F_{ji}=-F_{ij}\); the second uses the curvature commutator identity.
For any scalar smooth function \(\tau(s)\),

\[
 \sum_iD_i^{C+\tau S}(\dot\tau S_i)
 =\dot\tau\sum_iD_iS_i+\tau\dot\tau\sum_i[S_i,S_i]=0.
\]

Thus Gauss holds for (HC16). Expand the spatial curvature as
\(F+\tau B+\tau^2 C\) and its divergence as
\(S+\tau L+\tau^2Q+\tau^3N\). For \(\tau=s^2/2\), the temporal term
is \(-\ddot\tau S=-S\), giving (HC17) exactly. All derivatives and
products have support inside \(\operatorname{supp}C\); every coefficient
is smooth and bounded on that compact set. This proves support and
finite energy. At \(s=0\) the core has nonzero commutator curvature,
and continuity preserves it for a nonempty time interval. \(\square\)

This is a concrete correction with an exactly computed residual. The
residual in (HC17) has not been set to zero. On the equal-colour core,
\(C_j=e_j\), \(S_j=-8e_j\), and \(A_j=(1-4s^2)e_j\); direct insertion gives

\[
 D_A^\mu F_{\mu j}
 =(96s^2-384s^4+512s^6)e_j,
 \quad (L_j,Q_j,N_j)=(192,-1536,4096)e_j.
 \tag{HC20}
\]

The next Taylor coefficient is already fixed by this calculation:
\(a_j=s^2S_j/2+s^4L_j/24+\cdots\). Proving convergence or controlling
the full support boundary requires more than those Taylor coefficients.

\subsection*{7. Exact source-free evolution of every homogeneous core}

Restrict (HC1) to \(B_r(0)\). Its three components are the constants

\[
 c_1=\mu_{\mathbf i}(\alpha),\quad
 c_2=\mu_{\mathbf j}(\beta),\quad
 c_3=\mu_{\mathbf k}(\gamma).
 \tag{HC21}
\]

This restriction is a proved map from the full labelled fibre. Consider
the following separate spatially homogeneous receiving problem, retaining
the original \(c_i\):

\[
 \ddot V_j=\sum_i[V_i,[V_i,V_j]],\qquad
 V_j(0)=c_j,\quad\dot V_j(0)=0.
 \tag{HC22}
\]

Theorem 7.1. Equation (HC22) has a unique smooth solution for all
\(s\in\mathbb R\), depends smoothly on every initial \(c_i\), and is
\(\operatorname{Sp}(1)\)-equivariant. The four-dimensional field
\(A_0=0\), \(A_j=\varphi(V_j(s))\) solves every source-free Yang–Mills
equation on \(\mathbb R\times\mathbb R^3\), and on
\(\mathbb R\times T^3_{\ell_1,\ell_2,\ell_3}\) for any stated positive
side lengths. No division by a \(c_i\) occurs, so the construction applies
to all initial rank strata of \(\mathcal M_H\).

Proof. The right side of (HC22) is polynomial in nine real coordinates,
so a local solution follows from the contraction of its integral equation
on a sufficiently short interval: on a closed ball the polynomial vector
field and its derivative are bounded; choosing the interval so that the
derivative bound times its length is less than one gives existence and
uniqueness. Differentiating that integral equation gives smooth dependence
on the initial data and smoothness in time.

Let \(\langle\cdot,\cdot\rangle\) be the original quaternionic Euclidean
metric and define

\[
 \mathcal E=\frac12\sum_j|\dot V_j|^2
       +\frac12\sum_{i<j}|[V_i,V_j]|^2.
 \tag{HC23}
\]

The identity \(\langle[x,y],z\rangle=-\langle y,[x,z]\rangle\) gives

\[
 \frac{d}{ds}\frac12\sum_{i<j}|[V_i,V_j]|^2
 =-\sum_j\left\langle
       \sum_i[V_i,[V_i,V_j]],\dot V_j\right\rangle.
\]

Adding the derivative of the kinetic term and using (HC22) proves
\(\dot{\mathcal E}=0\). Consequently
\(|\dot V_j|\le\sqrt{2\mathcal E(0)}\) and
\(|V_j(s)|\le|c_j|+|s|\sqrt{2\mathcal E(0)}\).
On each finite time interval these bounds keep the integral equation in
a bounded region. Its polynomial derivative is bounded there, so the
local construction extends through every finite endpoint. This proves
global existence, uniqueness, and smooth dependence on each finite
time interval.

Set \(\mathcal G=\sum_j[V_j,\dot V_j]\). Its derivative is
\(\sum_{i,j}[V_j,[V_i,[V_i,V_j]]]\). For a pair \(i\ne j\), write
\(Z=[V_i,V_j]\). The two ordered contributions are
\([V_j,[V_i,Z]]-[V_i,[V_j,Z]]=[[V_j,V_i],Z]=[-Z,Z]=0\).
The diagonal terms vanish. Thus \(\dot{\mathcal G}=0\), and
\(\mathcal G(0)=0\) proves Gauss's law. The spatial Yang–Mills equation
for a homogeneous field is precisely (HC22), with its sign fixed by
\(D^0=-\partial_s\). Applying the exact Lie isomorphism \(\varphi\)
proves the matrix equation.

Simultaneous conjugation by \(u\in\operatorname{Sp}(1)\) preserves
the bracket and (HC22). Uniqueness therefore proves equivariance.
The coefficients are independent of \(x\), so they descend to every
stated spatial torus. \(\square\)

The physical energy density, with every trace and coupling factor, is

\[
 \varepsilon=\frac1{2g^2}\left(
 -2\sum_j\operatorname{tr}(F_{0j}^2)
 -2\sum_{i<j}\operatorname{tr}(F_{ij}^2)\right)
 =\frac4{g^2}\mathcal E.
 \tag{HC24}
\]

Its torus energy is \(\ell_1\ell_2\ell_3\varepsilon\). If
\(\mathcal E>0\), its energy on \(\mathbb R^3\) is infinite. The map
from (HC21) to (HC22) therefore solves the homogeneous core receiving
problem; it does not identify the full compactly supported profile with
a homogeneous field. The cutoff derivative terms of (HC14)–(HC18)
remain in the full problem.

Equivariance gives the same construction over \(D\) and over \(X\).
It gives a smooth map from their full rank-24 bundles into homogeneous
solution families, without selecting a nowhere-zero global section.
Pulling the base connection of Section 4 to this family gives a full
parameter-space connection. Its mixed curvatures are
\(\partial_A\varphi(V_i)+[b_A,\varphi(V_i)]\); they have not been
declared zero.

\subsection*{8. The equal-colour orbit and the retained source constants}

For the core \(c_j=e_j\), uniqueness and substitution give
\(V_j=f(s)e_j\), where

\[
 \ddot f+8f^3=0,\qquad f(0)=1,\quad\dot f(0)=0,
 \qquad \dot f^2+4f^4=4.
 \tag{HC25}
\]

The first integral follows by multiplication by \(2\dot f\).
Starting at \(f=1\), the acceleration is \(-8\), so \(f\) decreases;
on that branch \(\dot f=-2\sqrt{1-f^4}\). Its first zero occurs at

\[
 q=\frac12\int_0^1\frac{du}{\sqrt{1-u^4}}<\infty.
 \tag{HC26}
\]

The endpoint singularity is integrable since
\(1-u^4=(1-u)(1+u+u^2+u^3)\). Reflection at the turning points and
odd continuation through zeros, justified by uniqueness of the smooth
ODE, gives a periodic solution of period \(4q\). This is a complete
quadrature construction, without an imported elliptic-function identity.
Its Taylor coefficients begin
\(f(s)=1-4s^2+8s^4+O(s^6)\), from differentiating (HC25).

In the original matrix convention,

\[
 F_{0j}=2\dot f T_j,\quad
 (F_{12},F_{23},F_{31})=4f^2(T_3,T_1,T_2),\quad
 \varepsilon=\frac6{g^2}(\dot f^2+4f^4)=\frac{24}{g^2}.
 \tag{HC27}
\]

The magnetic density is \(48f^4\), the electric density is \(12\dot f^2\),
and their sum is \(48\). At the isolated zeros of \(f\), the magnetic
curvature vanishes while the electric part is nonzero. At \(s=0\), the
temporal contribution \(-2\ddot f T_j=16T_j\) cancels the exact static
source \(-16T_j\). This proves the cancellation in the original units;
in the physical time \(t\), the ODE is \(c^{-2}f_{tt}+8f^3=0\).

The homogeneous mechanism is established literature, not a novelty claim.
A. Tsapalis, E. P. Politis, X. N. Maintas, and F. K. Diakonos derive the
temporal-gauge homogeneous equations in
\href{https://arxiv.org/abs/1603.01858}{their original paper}, equations
(A15)–(A17) and (B8) in its TeX labels. Their source has Hermitian
\(t_a=\sigma_a/2\), physical field
\(\mathcal A_i^a=(m/g)\delta_i^a\Phi(ms)\), and
\(\Phi''+2\Phi^3=0\), in units where their time coordinate is a length.
The exact comparison to our anti-Hermitian connection is

\[
 A_i=-ig\mathcal A_i^a t_a=m\Phi(ms)T_i=2f(s)T_i,
 \qquad f(s)=\frac m2\Phi(ms).
 \tag{HC28}
\]

Then \(f''+8f^3=(m^3/2)(\Phi''+2\Phi^3)=0\), with
\(\Phi(0)=2/m\), \(\Phi'(0)=0\). The curvature comparison is
\(F=-ig\mathcal F^a t_a\), and their energy density
\(\sum_a(E_a^2+B_a^2)/2\) equals (HC24). Their metric is the negative
of ours; changing that overall sign leaves the source-free equation
unchanged. The value \(m>0\) remains explicit. It is not interpreted
as a quantum mass gap.

\subsection*{9. The static obstruction and the next boundary calculation}

Proposition 9.1. A smooth, compactly supported, time-independent
connection \(A_0=0\) on \(\mathbb R^{1,3}\) that solves the source-free
Yang–Mills equation has \(F_{ij}=0\).

Proof. Keep the magnetic energy

\[
 E(A)=\frac1{2g^2}\int_{\mathbb R^3}
       \sum_{i<j}\langle F_{ij},F_{ij}\rangle_{\mathfrak{su}(2)}\,d^3x,
 \qquad\langle U,V\rangle_{\mathfrak{su}(2)}=-2\operatorname{tr}(UV).
\]

For compactly supported variation \(a\), integration by parts gives

\[
 \delta E(A)[a]=-\frac1{g^2}\int\sum_j
    \left\langle\sum_iD_iF_{ij},a_j\right\rangle\,d^3x.
\]

The boundary term vanishes because both fields have compact support.
For the exact variation \(A_i^{(\lambda)}(x)=\lambda A_i(\lambda x)\),
\(\lambda>0\), one has
\(F_{ij}^{(\lambda)}(x)=\lambda^2F_{ij}(\lambda x)\), hence
\(E(A^{(\lambda)})=\lambda E(A)\). Its derivative at \(\lambda=1\)
uses \(a_i=A_i+x^k\partial_kA_i\), which has compact support. The
source-free equation makes this derivative zero; the exact scaling
makes it \(E(A)\). Positivity of the stated metric gives \(F=0\).
\(\square\)

Thus an interacting correction with the retained compact support cannot
remain static. Sections 6–8 have already attempted the next construction:
an exact Gauss-preserving compactly supported path with its complete
residual, and an exact global source-free evolution of the core. The
next calculation at this checkpoint was to solve (HC15) through the original cutoff region,
controlling the residual (HC17), spatial propagation, finite energy, and
parameter derivatives. The soft-period family must then be coupled to
this evolution with its full period matrix and monodromy. None of the
classical energies above is a spectral measure of a reconstructed
quantum Hamiltonian.

\subsection*{10. Sources, checks, and figure}

The higher-rung source and \href{https://github.com/KokunoYumeto/yang-mills-interacting-workbench/blob/main/yang-mills/continuations/20260930-s6-ns-moment-map-bridge/PROOF.md}{PROOF.md} supply the retained
topological class, representation, and original moment-map construction.
The new pullback formulas and evolution calculations are proved above.
Ichiro Yokota's
\href{https://arxiv.org/abs/0902.0431}{Exceptional Lie groups}, Theorems
1.16.2 and 2.7.1, supplies the human-source group conventions; its original
TeX was read at those proofs. The temporal-gauge source comparison uses
only Sections II–III of the Tsapalis–Politis–Maintas–Diakonos TeX,
including its equations (A1)–(A18), (B1)–(B4), and (B8). Both original
arXiv source archives are preserved locally; the reading ledger identifies
their byte hashes and records the unread sections.

The reproducible checker `checks/verify\_higher\_carrier\_evolution.py`
checks the polynomial source expansion, the homogeneous force and Gauss
identities, the energy cancellation, all equal-colour coefficients, and
the original curvature/trace factors. These checks support the algebra;
the bundle and ODE existence proofs are the arguments above.

\textit{The reproducible figure and inspected rendering are retained with the corresponding Markdown proof.}

Figure: the pullback square of Theorem 2.1 and the exact equal-colour
solution of (HC25)–(HC27). The plotted curve is a numerical rendering of
the proved ODE, not a numerical proof. Its labelled energies retain the
original factors. Reproducible source:
`figures/higher\_carrier\_evolution\_figure.py`.

\subsection*{Global compact-support Cauchy continuation}

The complete proof is \href{https://github.com/KokunoYumeto/yang-mills-interacting-workbench/blob/main/yang-mills/continuations/20260930-s6-ns-moment-map-bridge/COMPACT_SUPPORT_CAUCHY_EVOLUTION.md}{COMPACT\_SUPPORT\_CAUCHY\_EVOLUTION.md}.
Theorem 2.1 applies Oh's established global theorem to the original cutoff
data with zero electric field, checking every hypothesis and the metric,
trace, coupling and time conventions. Proposition 3.1 gives the support
bound \(R+|s|\), and Corollary 3.2 proves exact agreement with the
homogeneous solution on \(|x|+|s|<r\). The gauge-invariant core witness
in CE17 is at least 128, including at magnetic zeros.

Proposition 4.1 identifies the energy-zero inputs exactly with the original
rank-12 subbundle. CE22 retains all cutoff energy terms. Sections 5–6 prove
smooth parameter dependence, exact solution-map ranks and fibres, bundle
descent and every mixed curvature component. CE38–CE45 retain dilation,
support, energy and the ordered gauge-equivariant link-holonomy map.

The classical finite-energy Cauchy step is now supplied. The complete
period and monodromy coupling, physical gauge projection, actual Gram and
Hamiltonian kernels, quantum spectrum and continuum reconstruction remain
active. Classical energies tending to zero are not quantum spectral weight.
\end{document}
